% ---
% id: statistics/thm:geometry-validation
% kind: thm
% title: "Validity of the admissible geometries"
% summary: >-
%   An acceptance set whose data pass validation (a probability vector, or a
%   feasible linear system) is a non-empty, compact, admissible acceptance
%   set ([[meat-core/def:admissible-acceptance-set]]). For a box of
%   per-outcome bounds, validation means the lower sum to at most one and the
%   upper to at least one.
% prominence: main
% status: implemented
% uses: [def:taxonomy, def:matched-alphabet]
% mentions: [meat-core/def:admissible-acceptance-set]
% assumes: [A-ALPHABET]
% refs: []
% provenance: newly-derived
% implements: certiqs_key.statistics.taxonomy:validate_geometry
% lean:
% ---

The consuming engine assumes that its acceptance set is a closed convex subset
of the probability simplex. The alphabet is the finite, enumerated one of
Definition~\ref{def:matched-alphabet}, so the simplex is compact; a polytope cuts it
with finitely many closed half-spaces, and along the segment from the lower to
the upper bound vector of a canonical box the coordinate sum moves
continuously from the sum of the lower bounds to that of the upper bounds,
which gives the predicate in closed form.

\begin{theorem}\label{thm:geometry-validation}
Let $\Chatset$ and $n$ be a \tref{matched alphabet and round count}
(Definition~\ref{def:matched-alphabet}), and let $(\text{tag},\text{data})$
be an admissible geometry with denotation $\Gamma$.
\begin{enumerate}[label=(\roman*),leftmargin=2.2em]
\item If the validation predicate holds, $\Gamma$ is a non-empty, closed,
convex and compact subset of $\simplex{\Chatset}$. In particular
\xref{meat-core/def:admissible-acceptance-set}, that $\acSet$ is a closed convex subset of the
probability simplex on $\Chatset$, holds by construction for every admissible
geometry whose validation predicate holds.
\item For the tag $\mathrm{singleton}$, the validation predicate is necessary
and sufficient for $\Gamma$ to be a non-empty subset of $\simplex{\Chatset}$.
\item For the tag $\mathrm{polytope}$ in canonical box form with
$0\le\qlo\le\qhi$, the validation predicate, the condition
$\Gamma\neq\varnothing$ and the closed-form condition
\begin{equation}\label{eq:boxfeas}
  \sum_{\hat c\in\Chatset}q^{-}(\hat c)\ \le\ 1\ \le\
  \sum_{\hat c\in\Chatset}q^{+}(\hat c)
\end{equation}
are equivalent.
\end{enumerate}
\end{theorem}

\begin{proof}
\emph{(i).} $\simplex{\Chatset}$ is the intersection of the closed half-spaces
$\{q(\hat c)\ge0\}$ with the hyperplane $\{\sum_{\hat c}q(\hat c)=1\}$, hence
closed and convex; it is bounded, since $0\le q(\hat c)\le1$ for every
coordinate, hence compact ($\Chatset$ is finite, so
$\mathbb{R}^{\Chatset}$ is finite-dimensional and Heine--Borel applies). For
$\mathrm{polytope}$, $\{B\mathbf{q}\le\mathbf{b}\}$ is a finite intersection
of closed half-spaces, so closed and convex, and $\Gamma$ is the intersection
of two closed convex sets, hence closed and convex; being a closed subset of
the compact $\simplex{\Chatset}$, it is compact. For $\mathrm{singleton}$, a
one-point set is closed, convex and compact, and lies in $\simplex{\Chatset}$
exactly when the validation predicate holds. Non-emptiness follows from the
validation predicate in every case: for $\mathrm{singleton}$ by (ii); for a
general $\mathrm{polytope}$ because the predicate is rational feasibility of
$\{\mathbf{q}\ge0,\ \sum_{\hat c}q(\hat c)=1,\ B\mathbf{q}\le\mathbf{b}\}$ and
a feasible point is by \eqref{eq:polytope} an element of $\Gamma$; and in
canonical box form equivalently by the closed form (iii).

\emph{(ii).} $\{\mathbf{q}^{\ast}\}\subseteq\simplex{\Chatset}$ and
$\{\mathbf{q}^{\ast}\}\neq\varnothing$ hold if and only if
$\mathbf{q}^{\ast}\in\simplex{\Chatset}$, which is the predicate verbatim.

\emph{(iii).} We show that the predicate gives $\Gamma\neq\varnothing$, that
$\Gamma\neq\varnothing$ gives \eqref{eq:boxfeas}, and that
\eqref{eq:boxfeas} gives the predicate. By \eqref{eq:boxrows} the rows
$B\mathbf{q}\le\mathbf{b}$ read $\qlo\le\mathbf{q}\le\qhi$ componentwise.

\emph{Predicate $\Rightarrow$ non-emptiness.} A rational solution of the
system of the predicate is by \eqref{eq:polytope} an element of $\Gamma$.

\emph{Non-emptiness $\Rightarrow$ \eqref{eq:boxfeas}.} If
$\mathbf{q}\in\Gamma$ then $\qlo\le\mathbf{q}\le\qhi$ componentwise and
$\sum_{\hat c}q(\hat c)=1$; summing the two componentwise inequalities gives
\eqref{eq:boxfeas}.

\emph{\eqref{eq:boxfeas} $\Rightarrow$ predicate.} Put, for
$\theta\in[0,1]$, $\mathbf{q}(\theta)\coloneqq\qlo+\theta(\qhi-\qlo)$. Since
$\qlo\le\qhi$, every $\mathbf{q}(\theta)$ satisfies
$\qlo\le\mathbf{q}(\theta)\le\qhi$ and, using $\qlo\ge0$, also
$\mathbf{q}(\theta)\ge0$. Write $S^{-}\coloneqq\sum_{\hat c}q^{-}(\hat c)$ and
$S^{+}\coloneqq\sum_{\hat c}q^{+}(\hat c)$, so that $S^{-}\le1\le S^{+}$ by
\eqref{eq:boxfeas} and $\sum_{\hat c}q(\theta)(\hat c)=S^{-}+\theta(S^{+}-S^{-})$.
Put $\theta^{\ast}\coloneqq(1-S^{-})/(S^{+}-S^{-})$ if $S^{+}>S^{-}$, and
$\theta^{\ast}\coloneqq0$ if $S^{+}=S^{-}$, in which case $S^{-}=S^{+}=1$. In
both cases $\theta^{\ast}\in[0,1]$ and
$\sum_{\hat c}q(\theta^{\ast})(\hat c)=1$, and $\theta^{\ast}$ is rational
because $\qlo$ and $\qhi$ are (Definition~\ref{def:taxonomy}). So
$\mathbf{q}(\theta^{\ast})$ is a rational vector with
$\mathbf{q}(\theta^{\ast})\ge0$, $\sum_{\hat c}q(\theta^{\ast})(\hat c)=1$ and
$\qlo\le\mathbf{q}(\theta^{\ast})\le\qhi$, a rational solution of the system
of the predicate.
\end{proof}
